% Alcubierre's warp bubble (radius R = 1, wall parameter sigma = 4, speed v = 1)
% in a plane through the direction of travel (x) and a transverse direction
% (y). The bubble is at the origin and moves toward +x. With r = sqrt(x^2+y^2)
% and the shape function
%   f(r)  = [tanh(s(r+R)) - tanh(s(r-R))] / (2 tanh(s R)),
%   f'(r) = s [sech^2(s(r+R)) - sech^2(s(r-R))] / (2 tanh(s R)),
% (a) the expansion of the volume of space, theta = v (x/r) f'(r);
% (b) the energy density measured by observers at rest in the grid,
%     rho = -(v^2 / 32 pi) (y^2 / r^2) f'(r)^2, shown as -rho / max(-rho).
% (The sqrt argument is offset slightly: pgf's floating-point sqrt(0) fails.)
\documentclass[tikz,border=1pt]{standalone}
% Shared preamble for the standalone figures: the book's fonts, palette and
% drawing styles. Figures are drawn at their final printed size. A column is
% 3.35in (241pt) wide; a figure across both columns is 7in (505pt) wide.

\usepackage[T1]{fontenc}
\usepackage{newpxtext}
\usepackage[varbb]{newpxmath}
\usepackage[semibold,scale=0.95,lining]{sourcesanspro}
\usepackage{xcolor}
% Palette shared by the book and by the standalone figures in figures/src.
% One main hue (blue), one accent (orange), and a few roles used sparingly.

\definecolor{seblue}{HTML}{1D5B8F}    % headings, worked examples, main curves
\definecolor{seblueDark}{HTML}{123E63}
\definecolor{seorange}{HTML}{DC7A26}  % thought experiments, light and light cones
\definecolor{seteal}{HTML}{1F8577}    % checks, travellers and ships
\definecolor{sered}{HTML}{C0392B}     % negative energy, tension, violations
\definecolor{seviolet}{HTML}{5E548E}  % going deeper
\definecolor{sesand}{HTML}{8C6A3F}    % history
\definecolor{segray}{HTML}{6B7480}    % axes, guides, secondary labels
\definecolor{selight}{HTML}{D5DAE1}   % grids and faint rules

\usepackage{siunitx}
\usepackage{fontawesome5}
\usepackage{tikz}
\usetikzlibrary{arrows.meta,calc,positioning,decorations.pathreplacing,
  decorations.markings,shapes.geometric,fadings,shadings,patterns,3d,
  intersections,backgrounds}
\usepackage{pgfplots}
\pgfplotsset{compat=1.18}
\usepgfplotslibrary{fillbetween,groupplots,colormaps}

\tikzset{
  >={Latex[length=2.1mm,width=1.4mm]},
  every picture/.style={line cap=round, line join=round},
  lbl/.style={font=\sffamily\footnotesize, text=black!85},
  smalllbl/.style={font=\sffamily\scriptsize, text=black!75},
  guide/.style={segray, line width=0.4pt},
  axisline/.style={segray, line width=0.5pt, -{Latex[length=1.7mm,width=1.1mm]}},
  light/.style={seorange, line width=0.9pt},
  lightdash/.style={seorange, line width=0.9pt, dash pattern=on 3pt off 2pt},
  ship/.style={seteal, line width=1.4pt},
  cone/.style={fill=seorange!22, draw=seorange!85!black, line width=0.35pt},
}

\pgfplotsset{
  se axis/.style={
    axis lines=left,
    axis line style={segray, line width=0.5pt, -{Latex[length=1.6mm,width=1.1mm]}},
    tick style={segray, line width=0.4pt},
    major tick length=2.5pt,
    tick label style={font=\sffamily\scriptsize, text=black!75},
    label style={font=\sffamily\footnotesize, text=black!85},
    every axis plot/.append style={line width=1.1pt},
    legend style={font=\sffamily\scriptsize, draw=none, fill=none},
    scaled ticks=false,
    xticklabel={\sansnum{\tick}}, yticklabel={\sansnum{\tick}},
  },
}

% Numbers in figures are set in the sans-serif text font, with a true minus.
\sisetup{mode=text, reset-text-family=false, reset-text-series=false,
  reset-text-shape=false, drop-zero-decimal, group-minimum-digits=5}
% Tick values arrive as long decimals (1.5000000000); parse them first.
\newcommand{\sansnum}[1]{\begingroup\pgfkeys{/pgf/fpu=false}%
  \pgfmathparse{1*(#1)}\num{\pgfmathresult}\endgroup}

\begin{document}
\begin{tikzpicture}[
  declare function={
    rr(\x,\y)=sqrt(\x*\x+\y*\y+1e-10);
    sechsq(\a)=1/(cosh(\a)^2);
    fp(\r)=4*(sechsq(4*(\r+1)) - sechsq(4*(\r-1)))/(2*tanh(4));
  },
]
% ---------------------------------------------------------------- (a)
\begin{axis}[
  name=A, hide axis, view={-8}{30}, clip=false,
  width=318pt, height=230pt,
  xmin=-2.6, xmax=2.6, ymin=-2.6, ymax=2.6, zmin=-2.2, zmax=2.2,
  colormap={diverge}{rgb255=(29,91,143) rgb255=(96,146,192) rgb255=(172,200,225)
    rgb255=(236,238,240) rgb255=(242,196,152) rgb255=(230,150,82) rgb255=(214,112,32)},
  point meta min=-2.1, point meta max=2.1,
]
\addplot3[surf, shader=faceted interp, faceted color=black!18, line width=0.06pt,
  z buffer=sort, samples=73, samples y=73, domain=-2.6:2.6, y domain=-2.6:2.6]
  {x/rr(x,y)*fp(rr(x,y))};
\coordinate (behind) at (axis cs:-1.0,0,1.95);
\coordinate (ahead)  at (axis cs:1.15,0,-1.2);
\coordinate (arrowstart) at (axis cs:-1.3,-3.2,-1.2);
\coordinate (arrowend)   at (axis cs:1.5,-3.2,-1.2);
\end{axis}
\draw[guide] (behind) -- ++(-26pt,10pt) node[lbl, anchor=east, text=seorange!85!black,
  align=right] {space expands\\behind the bubble};
\draw[guide] (ahead) -- ++(30pt,-14pt) node[lbl, anchor=west, text=seblue, align=left]
  {space contracts\\ahead of it};
\draw[->, black!70, line width=0.8pt] (arrowstart) -- (arrowend)
  node[pos=0.5, below=1pt, smalllbl] {direction of travel};
\node[font=\sffamily\bfseries\small, text=black!85, anchor=north west] at (A.north west) {(a)};

% ---------------------------------------------------------------- (b)
\begin{axis}[
  name=B, at={($(A.east)+(14pt,8pt)$)}, anchor=west,
  hide axis, view={0}{90}, axis equal image,
  width=200pt, height=200pt,
  xmin=-1.8, xmax=1.8, ymin=-1.8, ymax=1.8,
  colormap={neg}{rgb255=(255,255,255) rgb255=(248,222,214) rgb255=(238,170,150)
    rgb255=(212,98,78) rgb255=(168,50,36) rgb255=(120,24,18)},
  point meta min=0, point meta max=1,
  colorbar horizontal,
  colorbar style={
    width=110pt, height=5pt, at={(0.5,-0.04)}, anchor=north,
    xtick={0,1}, xticklabels={zero, most negative},
    tick label style={font=\sffamily\scriptsize, text=black!75},
    xticklabel style={yshift=1pt}, major tick length=0pt,
    axis line style={draw=none},
    xlabel={energy density demanded}, xlabel style={font=\sffamily\scriptsize,
      text=black!75, yshift=3pt},
  },
]
\addplot3[surf, shader=interp, samples=121, samples y=121,
  domain=-1.8:1.8, y domain=-1.8:1.8]
  {(y*y/(rr(x,y)^2))*fp(rr(x,y))^2/4.02};
\draw[black!45, line width=0.45pt, dash pattern=on 2pt off 1.5pt]
  (axis cs:0,0) circle[radius=1];
\fill[seteal] (axis cs:-0.3,-0.075) rectangle (axis cs:0.3,0.075);
\draw[seteal, line width=0.6pt, -{Latex[length=1.5mm,width=1.2mm]}]
  (axis cs:0.36,0) -- (axis cs:0.72,0);
\node[smalllbl, text=black!80] at (axis cs:0,0.3) {ship};
\coordinate (wall) at (axis cs:0.62,-0.87);
\end{axis}
\draw[guide] (wall) -- ++(30pt,-18pt) node[smalllbl, anchor=west, align=left] {bubble\\wall};
\node[font=\sffamily\bfseries\small, text=black!85, anchor=north west] at ($(B.north west)+(-6pt,-2pt)$) {(b)};
\end{tikzpicture}
\end{document}
